\RequirePackage{fix-cm}
\documentclass[10pt]{article}
\usepackage[margin=0.75in]{geometry}
\usepackage{amsmath,amssymb,amsthm,mathtools,enumitem}
\usepackage{fvextra}
\setlist{nosep,leftmargin=1.5em}
\usepackage{microtype}
\newtheoremstyle{tplain}{1.5pt}{0.5pt}{\itshape}{}{\bfseries}{.}{0.5em}{}
\newtheoremstyle{tdef}{1.5pt}{0.5pt}{\normalfont}{}{\bfseries}{.}{0.5em}{}
\theoremstyle{tplain}\newtheorem{theorem}{Theorem}\newtheorem{lemma}[theorem]{Lemma}\newtheorem{proposition}[theorem]{Proposition}\newtheorem{corollary}[theorem]{Corollary}
\theoremstyle{tdef}\newtheorem{definition}[theorem]{Definition}
\makeatletter\renewenvironment{proof}[1][\proofname]{\par\pushQED{\qed}\normalfont\topsep1pt\partopsep0pt\trivlist\item[\hskip\labelsep\itshape#1\@addpunct{.}]\ignorespaces}{\popQED\endtrivlist\@endpefalse}\makeatother
\AtBeginDocument{\setlength\abovedisplayskip{3pt}\setlength\belowdisplayskip{3pt}\setlength\abovedisplayshortskip{1pt}\setlength\belowdisplayshortskip{1pt}}
\newcommand{\inp}{\mathrm{in}}\newcommand{\outp}{\mathrm{out}}\newcommand{\ports}{\mathrm{ports}}
\newcommand{\priv}{\mathrm{priv}}\newcommand{\edge}{\mathrm{edge}}
\newcommand{\TRUE}{\mathsf{true}}\newcommand{\FALSE}{\mathsf{false}}
\newcommand{\hid}[1]{\overline{#1}}
\newcommand{\df}[1]{[\![#1]\!]}
\setlength{\parskip}{0.5pt}\setlength{\parindent}{0pt}
\makeatletter\renewcommand\section{\@startsection{section}{1}{\z@}{3pt}{1pt}{\normalfont\bfseries}}\makeatother
\begin{document}
\begin{center}{\large\bfseries Computing local delegations: the complete procedure}\\[1pt] 18 October 2026\end{center}

\section{The problem}
The problem, in the author's words:\vspace{-2pt}
\begin{quote}\fontsize{8}{9.5}\selectfont
We have a system which is a graph of components where each component has labeled/typed inputs and outputs (``typed'' means the variables are distinct and each is assigned to a known component; there are no numeric range bounds). We have a global query contract written in PolySTL (polytopic signal temporal logic: STL whose predicates are convex polytopes over the state variables, $\mu(x):=Ax\le b$, where $x$ is the vector of ALL system variables). We need to ALGEBRAICALLY DECOMPOSE the global contract into obligations for the components. The delegate contract handed to a component may ONLY reference the variables that component has available, i.e.\ its own inputs and outputs. Only some components will receive a (non-trivial) delegated contract: not every global contract needs every component. All that is available is: (1) the architectural graph (components, typed/labelled ports, who reads/writes which variable), and (2) the top-level query contract. No component models.
\end{quote}\vspace{-3pt}
The composition of the delegated contracts must refine the query, and the algebra may use only contract operations: refinement, composition, quotient and projection.

\textbf{The solution in plain words.} Relay: Radio reads $r$ and sends $m$ to Display, which outputs $d$; the query is ``if $|r|\le1$ then $|d-r|\le\frac12$''. Refinement by the composition amounts to two implications. Contracts no component can meet satisfy them, so each contract must also be met by a component computing its outputs from the inputs received so far. None is best; the procedure picks one. \emph{Forwarding:} each component copies onto its outgoing variables what the next cannot see ($m=r$). \emph{Choice:} one linear-arithmetic sentence says that for every input each component in turn has an output making the query true; eliminating its quantifiers gives each output a rule, the first of fixed candidates that works ($d=r$, in Display's input $d=m$). \emph{Enlargement:} quotient and projection enlarge each contract until none can be enlarged alone. If the query mentions only ports of relevant components forming a forwarding chain, this is complete and a false sentence proves none exists; otherwise it can refute, and for $H=0$ search without a halting guarantee (Section~\ref{sec:scope}).

\section{Objects}\label{sec:obj}
\textbf{Assumptions and variables.} $V$ is a finite set of real-valued variables without range bounds. Each is written by exactly one component or by the environment (an \emph{environment variable}). Component $i$ reads $\inp(i)\subseteq V$ and writes $\outp(i)\subseteq V$; $\ports(i):=\inp(i)\cup\outp(i)$. The \emph{architecture graph} has an edge $u\to w$ when $w$ reads a variable that $u$ writes. We assume: (i) discrete time $t\in\{0,\dots,H\}$; (ii) no directed cycle; (iii) zero delay: an output at time $t$ may depend on inputs at time $t$. The \emph{timed variable} $v(t)$ is the value of $v$ at time $t$. A \emph{trace} assigns a real number to every timed variable; a \emph{history} of $W\subseteq V$ assigns one to every $v(t)$ with $v\in W$; $\hid W$ is the set of timed variables $v(t)$ with $v\notin W$.

\textbf{The formula language $\mathcal L$.} An \emph{atom} is $\sum_{v\in V}a_v\,v(t)\le b$, rational $a_v,b$, one time $t$: a half-space over the vector of all variables at time $t$. $\TRUE$ and $\FALSE$ are the atoms $0\le0$ and $0\le-1$. Formulas of $\mathcal L$ are built from atoms with $\neg,\wedge,\vee,\to$; a strict inequality is a negated atom; $e=c$ abbreviates $e\le c\wedge -e\le -c$, and $|e|\le c$ abbreviates $e\le c\wedge -e\le c$. A formula is identified with the set of traces satisfying it. It is \emph{over} $W$ if it mentions only $v(t)$ with $v\in W$; its truth then depends only on the history of $W$. $\varphi\models\psi$: every trace satisfying $\varphi$ satisfies $\psi$; $\equiv$: both ways. $\exists v(t)$ ranges over $\mathbb R$; $\exists\hid W$ quantifies all of $\hid W$. $\varphi[y:=e]$ replaces the timed variable $y$ by the expression $e$. A formula with quantifiers is not in $\mathcal L$, but by Lemma~\ref{lem:elim}(a) it is equivalent to one; below it always stands for that formula.

\textbf{Bounded PolySTL.} $\psi::=\mu\mid\neg\psi\mid\psi\wedge\psi\mid\psi\,\mathcal U_{[a,b]}\psi$, with $\mu$ a polytope atom $a\cdot x\le b$ over the vector $x$ of all variables. Its meaning at time $t$ is $[\psi]_t\in\mathcal L$: $[\mu]_t$ replaces each $v$ by $v(t)$; $[\cdot]_t$ commutes with $\neg,\wedge$; $[\psi_1\mathcal U_{[a,b]}\psi_2]_t:=\bigvee_{s=t+a}^{t+b}\big([\psi_2]_s\wedge\bigwedge_{t\le s'<s}[\psi_1]_{s'}\big)$; $\square_{[a,b]}$, $\Diamond_{[a,b]}$ give $\bigwedge$, $\bigvee$ over $s\in[t+a,t+b]$. ``Bounded'': all times so referenced from $t=0$ are $\le H$; a trace satisfies $\psi$ if it satisfies $[\psi]_0$. \emph{So $\mathcal L$ and bounded PolySTL define the same sets of traces}: $[\psi]_0\in\mathcal L$, and an atom at time $t$ is $[\Diamond_{[t,t]}\mu]_0$.

\textbf{Contracts.} A \emph{contract} for component $i$ is a pair $(A_i,G_i)$ of formulas over $\ports(i)$: assumption and guarantee. The \emph{query} is $P=(A_P,G_P)$, formulas over $V$. Its \emph{implementation set} is $S:=\neg A\vee G$, the traces on which it is not violated; an implementation is any set of traces inside $S$. It is \emph{saturated} if $G\equiv S$ (the guarantee allows every trace violating the assumption). Replacing $G$ by $S$ changes neither implementations nor assumption, so all contracts below are saturated, written $(A,S)$ with $\neg A\models S$; the query is $(A_P,S_P)$, $S_P:=\neg A_P\vee G_P$. The operations:
\[\begin{array}{ll}
\text{refinement:} & (A,S)\preccurlyeq(A',S')\ \text{iff}\ A'\models A\ \text{and}\ S\models S';\qquad\text{projection:}\ \ \rho_W(A,S):=\big(\exists\hid W\,A,\ \forall\hid W\,S\big);\\
\text{composition:} & (A_1,S_1)\otimes(A_2,S_2):=\big((A_1\wedge A_2)\vee\neg(S_1\wedge S_2),\ S_1\wedge S_2\big);\\
\text{quotient:} & (A,S)/(A',S'):=\big(A\wedge S',\ (S\wedge A')\vee\neg(A\wedge S')\big).
\end{array}\]
In words: $C$ refines $C'$ if it accepts every context (environment behaviour) $C'$ accepts and promises at least as much; composition puts two components side by side, promising both guarantees and needing both assumptions unless a promise is broken; $C/C'$ is what a component must meet so that, next to one meeting $C'$, the pair meets $C$; projection keeps what can be said over $W$. Each result is saturated (e.g.\ $\neg\exists\hid WA\equiv\forall\hid W\neg A\models\forall\hid WS$). By induction, $\bigotimes_{i}(A_i,S_i)=\big(\bigwedge_iA_i\vee\neg\bigwedge_iS_i,\ \bigwedge_iS_i\big)$: in the step, $\big((\bigwedge_{i<n}A_i\vee\neg\bigwedge_{i<n}S_i)\wedge A_n\big)\vee\neg\bigwedge_{i\le n}S_i\equiv\bigwedge_{i\le n}A_i\vee\neg\bigwedge_{i\le n}S_i$ since $\neg\bigwedge_{i<n}S_i\wedge A_n\models\neg\bigwedge_{i\le n}S_i$.
\begin{lemma}[adjunction, localisation]\label{lem:adj}
(a) $C_1\otimes C_2\preccurlyeq C$ iff $C_1\preccurlyeq C/C_2$. (b) $\rho_W(C)$ is the $\preccurlyeq$-greatest contract over $W$ that refines $C$.
\end{lemma}
\begin{proof}
(a) With $C=(A,S)$, both sides say $A\wedge S_1\wedge S_2\models A_1\wedge A_2\wedge S$ (left: use $\neg A\models S$; right: $\neg A_1\models S_1$). (b) For $(A',S')$ over $W$: $A\models A'$ iff $\exists\hid WA\models A'$, and $S'\models S$ iff $S'\models\forall\hid WS$, because $A',S'$ do not depend on $\hid W$.
\end{proof}

\section{What ``valid'' means}
A \emph{delegation} is a family of contracts $C_i=(A_i,S_i)$, $C_i$ over $\ports(i)$, one per component. It is \emph{valid} if $\bigotimes_iC_i\preccurlyeq P$.
\begin{proposition}[validity conditions]\label{prop:valid}
A delegation is valid iff (V1) $A_P\wedge\bigwedge_iS_i\models G_P$ and (V2) $A_P\wedge\bigwedge_iS_i\models A_j$ ($\forall j$).
\end{proposition}
\begin{proof}
By the formula for $\bigotimes$, validity means (i) $A_P\models\bigwedge_iA_i\vee\neg\bigwedge_iS_i$ and (ii) $\bigwedge_iS_i\models\neg A_P\vee G_P$. (ii) says every trace of $A_P\wedge\bigwedge_iS_i$ satisfies $G_P$: (V1). (i) says every trace of $A_P$ satisfying $\bigwedge_iS_i$ satisfies $\bigwedge_iA_i$: (V2).
\end{proof}
\begin{proposition}[validity alone is empty]\label{prop:empty}
(a) One component $(\TRUE,\FALSE)$, all others $(\TRUE,\TRUE)$, is valid. (b) If a rational trace $z$ satisfies $A_P\wedge G_P$ and all variables of $G_P$ are ports, $C_i:=\big(\TRUE,\bigwedge_{v\in\ports(i),t\le H}v(t)=z(v(t))\big)$ is valid.
\end{proposition}
(Proof: (a) $\bigwedge_iS_i\equiv\FALSE$; (b) a trace of $\bigwedge_iS_i$ agrees with $z$ on the variables of $G_P$.) Both are useless: in (a) the component given $(\TRUE,\FALSE)$ has no implementation; in (b) each component with an input forbids all input histories but one. A component must accept whatever arrives at its inputs, which depends on which ports are inputs, a distinction no operation uses:
\begin{definition}
$C_i$ is \emph{receptive} (R) if for every history of $\inp(i)$ some history of $\outp(i)$ gives a trace in $S_i$. A \emph{causal strategy} for $i$ is a family of functions $f_{i,t}$, $t\le H$, from the values of $\inp(i)$ at times $\le t$ to values of $\outp(i)$ at time $t$. $C_i$ is \emph{causal} (K) if some causal strategy produces, from every history of $\inp(i)$, a trace in $S_i$. A delegation is \emph{acceptable} if it satisfies (V1), (V2), and (R), (K) for every component. (K) implies (R); (R), (K) survive enlarging $S_i$; (V1), (V2) survive shrinking it.
\end{definition}

\begin{proposition}[no greatest solution]\label{prop:nogreatest}
In the \emph{relay} (Radio reads the environment variable $r$, writes $m$; Display reads $m$, writes $d$; $H=0$, $v$ for $v(0)$; $A_P=|r|\le1$, $G_P=|d-r|\le\frac12$), $D$: Radio $(\TRUE,m=r)$, Display $(\TRUE,d=m)$ and $D'$: Radio $(\TRUE,m=-r)$, Display $(\TRUE,d=-m)$ are acceptable, but no valid delegation has implementation sets containing those of both.
\end{proposition}
\begin{proof}
In $D$ and $D'$, $d=r$: (V1); (V2) is trivial; each output is a function of the input at the same time: (K), (R). A delegation containing both has $S_{\mathrm{Radio}}\supseteq(m=r\vee m=-r)$, $S_{\mathrm{Display}}\supseteq(d=m\vee d=-m)$; the trace $r=1$, $m=-1$, $d=-1$ satisfies $A_P$ and both, but $|d-r|=2$.
\end{proof}
\textbf{What the algebra alone gives.} By Lemma~\ref{lem:adj}, if all contracts but $C_i$ are fixed, then $\bigotimes_jC_j\preccurlyeq P$ iff $C_i\preccurlyeq\rho_{\ports(i)}\big(P/\bigotimes_{j\ne i}C_j\big)$: one quotient and one projection. With all others $(\TRUE,\TRUE)$, $P/(\TRUE,\TRUE)=(A_P,S_P\vee\neg A_P)=P$: the relay gets Radio $\rho_{\{r,m\}}(P)=(|r|\le1,|r|>1)$ and Display $\rho_{\{m,d\}}(P)=(\TRUE,\FALSE)$, neither receptive. No acceptable relay delegation comes from $P$ and constant contracts by the operations: these introduce no variable and $m$ is not in $P$, so Radio's set is over $\{r\}$, Display's over $\{d\}$; by (R) the first holds for all $r$, the second contains some $d_0$, and (V1) would need $|d_0-r|\le\frac12$ on $[-1,1]$. So the algebra relates each contract to the others but supplies no starting point, and by Proposition~\ref{prop:nogreatest} there is no best delegation for it to compute. Steps 2 to 7 therefore \emph{select} one (the only part that is not a contract operation, checked by Proposition~\ref{prop:valid}); steps 8, 9 use only quotient and projection.

\section{The computational core: elimination with witnesses}\label{sec:qe}
\textbf{Idea.} With all timed variables but $y$ fixed, each atom containing $y$ compares $y$ with one number, its bound; between consecutive bounds no atom changes truth, so $\varphi$ is constant there. It suffices to test each bound, a point between each two, points beyond the outermost, and one point if there is none; the candidates supply these for every arrangement.

Let $\varphi\in\mathcal L$ and $y=w(t)$. Each atom $\sum_va_v\,v(s)\le b$ of $\varphi$ with $a_w\ne0$ (so $s=t$) has the \emph{bound} $\big(b-\sum_{v\ne w}a_v\,v(t)\big)/a_w$, the value of $y$ at which it holds with equality; let $b_1,\dots,b_n$ be the distinct bounds. The \emph{candidates} $T$ are five tiers, in this order: (1) the midpoints $\frac12(b_a+b_{a'})$ of all pairs of distinct bounds, (2) the bounds $b_a$, (3) the bounds plus one, (4) the bounds minus one, (5) the constant $0$. Within a tier: first those whose value is a constant, by increasing absolute value, a negative constant before the positive one of the same absolute value; then the non-constant ones, in increasing character-code (ASCII) order of their canonical printing (terms $c$\texttt{*}$v$ in alphabetical order of $v$, then the constant; reduced fractions; joined by \texttt{' + '}; e.g.\ \texttt{'1*r\_0 + -3/8'} $<$ \texttt{'1*r\_0 + 3/8'} $<$ \texttt{'2*r\_0 + 0'}). A candidate equal to one met earlier is dropped. For $c\in T$ let $\pi_c:=\varphi[y:=c]\wedge\bigwedge_{c'\text{ before }c}\neg\varphi[y:=c']$ (the \emph{region} of $c$), and $\pi_0^\ast:=\bigwedge_{c\in T}\neg\varphi[y:=c]$. The \emph{witness} $s_y$ is the rule ``on $\pi_c$ take $c$; on $\pi_0^\ast$ take $0$'': given the other values, the first candidate at which $\varphi$ is true. Its \emph{defining formula} is $\df{y=s_y}:=\bigvee_{c}(\pi_c\wedge y=c)\vee(\pi^\ast_0\wedge y=0)$, true exactly when $y$ equals the value of $s_y$, and $\varphi[y:=s_y]:=\bigvee_c(\pi_c\wedge\varphi[y:=c])\vee(\pi^\ast_0\wedge\varphi[y:=0])$.
\begin{lemma}[elimination with witnesses]\label{lem:elim}
(a) $\exists y\,\varphi\equiv\bigvee_{c\in T}\varphi[y:=c]$ and $\forall y\,\varphi\equiv\bigwedge_{c\in T}\varphi[y:=c]$. (b) $\exists y\,\varphi\models\varphi[y:=s_y]$. (c) Each value $c$ of $s_y$ is affine at time $t$; $\varphi[y:=c]$, the regions and $\df{y=s_y}$ are in $\mathcal L$ (one time per atom).
\end{lemma}
\begin{proof}
With the other values fixed, the bound values $\beta_1<\dots<\beta_r$ cut $\mathbb R$ into cells (points $\beta_l$, open gaps, two rays; $\mathbb R$ if $r=0$), on each of which $\varphi$ is constant. Each cell contains a candidate: $\beta_l$; $\frac12(\beta_l+\beta_{l+1})$; $\beta_1-1$, $\beta_r+1$; $0$ if $r=0$. This gives (a) for $\exists$ and (b); for $\forall$ use $\neg\varphi$ (same atoms). (c): atoms with $y=w(t)$ mention only time $t$, so bounds are affine at time $t$.
\end{proof}
\textbf{Fourier--Motzkin} (for hand computation; witnesses always come from (b)). In a \emph{case}, a conjunction of atoms and negated atoms, each bound is a lower bound $\ell$ or an upper bound $u$ by the sign of $a_w$, and $\exists y\big(\gamma\wedge\bigwedge_\ell\ell\lhd_\ell y\wedge\bigwedge_uy\lhd_uu\big)\equiv\gamma\wedge\bigwedge_{\ell,u}\ell\lhd_{\ell u}u$ for $\gamma$ free of $y$, with $\lhd_{\ell u}$ strict iff either bound is: the admissible $y$ form an interval.


\section{Forwarding}\label{sec:fwd}
\textbf{Which components matter.} $N$ is the set of components that write a variable of $G_P$, together with every component from which a directed path leads to one of them. Every other component gets $(\TRUE,\TRUE)$: nothing it does reaches $G_P$. A writer of an input of a component in $N$ is in $N$. \emph{(Scope)}: every variable of $A_P$ or $G_P$ is a port of some component of $N$. A \emph{topological ordering} of $N$ is a list $i_1,\dots,i_p$ of its components in which each comes after every component whose outputs it reads. For $k<p$: $\priv_k:=\ports(i_k)\setminus\ports(i_{k+1})$, what $i_k$ sees and $i_{k+1}$ does not; $\edge_k:=\outp(i_k)\cap\inp(i_{k+1})$, the \emph{edge variables} from $i_k$ to $i_{k+1}$.

\textbf{Why one sentence with one quantifier order.} The exact question asks for one causal strategy per component, each reading only its own inputs. In a sentence $\forall\exists\forall\exists\cdots$ each existential choice may depend on everything to its left, which can grant values never received and make the sentence true although none exists (Section~\ref{sec:scope}). Forwarding removes the excess: if each component is followed by one it feeds (E), along an edge with at least as many variables as the first sees and the second does not (Cap), then $i_k$ copies those values onto the edge, and $i_{k+1}$ sees as copies all that $i_1,\dots,i_k$ saw (Lemma~\ref{lem:carry}). Knowledge is \emph{nested} along the list, so at one time one quantifier order grants each component only what it can know. As the sentence quantifies by time first, $i_1$'s choice at time 1 follows $i_2$'s inputs at time 0; (Info) stops $i_2$ reading the environment, which $i_1$ would be granted unseen. Earlier outputs of later components are granted too; under (Info) they are rules in values the chooser knows, which step 6 substitutes. Edge values are copies, not choices, so they are not quantified, and (Pure) keeps them out of the query. (In the relay Radio's only output is its edge $m$, so Radio chooses nothing.)
\begin{definition}[nested ordering]
A topological ordering is \emph{nested} if for every $k<p$: (E) $i_k\to i_{k+1}$ is an edge; (Cap) $|\priv_k|\le|\edge_k|$; (Pure) no edge variable occurs in $A_P$ or $G_P$; (Info) $H=0$, or every variable read by $i_k$ with $k\ge2$ is written by some $i_j$, $j<k$.
\end{definition}
(E) follows from (Cap): $i_k\in N$ writes a variable, in $\priv_k$ if $i_{k+1}$ reads none. The proofs use (Info) only as: for $H>0$, no $i_k$, $k\ge2$, reads an environment variable no earlier component reads. Fix injections $\sigma_k:\priv_k\to\edge_k$ pairing the two alphabetically sorted lists. The \emph{forwarding formula} is $F_k:=\bigwedge_{t\le H}\bigwedge_{v\in\priv_k}\sigma_k(v)(t)=v(t)$, over $\ports(i_k)$; $F_p:=\TRUE$; other edge variables are unconstrained (their values need have no physical meaning).
\begin{lemma}[carriers]\label{lem:carry}
Let $j<k$ be largest with $v\in\ports(i_j)$. Put $w_j:=v$ and, for $j\le l<k$, $w_{l+1}:=w_l$ if $w_l\in\ports(i_{l+1})$, else $w_{l+1}:=\sigma_l(w_l)$. Then $\kappa_k(v):=w_k\in\inp(i_k)$ and $F_j\wedge\dots\wedge F_{k-1}\models\kappa_k(v)(t)=v(t)$ for all $t$.
\end{lemma}
\begin{proof}
By induction $w_l\in\ports(i_l)$; if $w_l\notin\ports(i_{l+1})$ then $w_l\in\priv_l$, $\sigma_l$ is defined, and $F_l\models w_{l+1}(t)=w_l(t)$. So $w_k\in\ports(i_{k-1})\cap\ports(i_k)$ or $w_k\in\edge_{k-1}$; an output of $i_k$ is not a port of the earlier $i_{k-1}$, so $w_k\in\inp(i_k)$.
\end{proof}
So every value an earlier component saw arrives as a copy on an input $\kappa_k(v)$ of $i_k$. The \emph{decoding} $\delta_k(\varphi)$ replaces every $v(s)$ with $v$ a port of some $i_j$, $j<k$, and $v\notin\ports(i_k)$, by $\kappa_k(v)(s)$. On traces satisfying $F_1,\dots,F_{k-1}$, $\varphi$ and $\delta_k(\varphi)$ have the same truth value; $\delta_k$ keeps one time per atom.

\section{The procedure}\label{sec:proc}
Input: the architecture graph and $P$. Output: an acceptable delegation, ``none exists'', or ``does not apply''.\vspace{-1pt}
\begin{enumerate}
\item \emph{Relevance.} Write $A_P,G_P$ in $\mathcal L$. Compute $N$; outside $N$: $(\TRUE,\TRUE)$. If (Scope) fails, stop: does not apply.
\item \emph{Ordering.} Enumerate the permutations of $N$, sorted by name, in lexicographic order; take the first that is topological and nested. If none, stop: does not apply (Section~\ref{sec:scope} still offers a refutation test).
\item \emph{Sentence.} $E_k(t)$ is the set of timed variables $v(t)$ with $v\in\inp(i_k)$ not a port of any $i_j$, $j<k$ (\emph{fresh inputs}; a writer of $v$ would be in $N$ before $i_k$, so $v$ is an environment variable; under (Info) $E_k(t)=\emptyset$ for $k\ge2$ if $H>0$). $Y_k(t)$: the $v(t)$ with $v\in\outp(i_k)\setminus\edge_k$ (\emph{terminal outputs}: not on the edge to $i_{k+1}$; $\edge_p:=\emptyset$). The \emph{prefix} lists these timed variables by $t$, then $k$, each set in port declaration order. Let
\[\Phi:=\forall E_1(0)\,\exists Y_1(0)\,\forall E_2(0)\,\exists Y_2(0)\cdots\forall E_p(H)\,\exists Y_p(H)\ (A_P\to G_P).\]
$\Phi$ is a sentence: a port of $i_k$ is an output (in $Y_k$, or in $\edge_k$, excluded by (Pure)), a fresh input, or a port of an earlier component, so by (Scope) all of $A_P\to G_P$ is quantified.
\item \emph{Elimination (Lemma~\ref{lem:elim}).} $\psi:=A_P\to G_P$. For the pairs $(t,k)$ in reverse order: for each $y\in Y_k(t)$, last first, record the witness $s_y$ of $y$ in $\psi$ and replace $\psi$ by $\exists y\,\psi$; then eliminate $\forall$ for each timed variable of $E_k(t)$. Later variables are already eliminated, so $s_y$ mentions only earlier ones. The final $\psi$ is the truth value of $\Phi$.
\item \emph{Decision.} If $\Phi$ is false, stop: none exists (Theorem~\ref{thm:nec}).
\item \emph{Substitution.} For $y\in Y_k(t)$, take $s_y$ as its list of pairs (region, value). If $z(s)$ in it is a terminal output of some $i_j$, $j>k$, then $s<t$ ($z(s)$ precedes $y$), so $z(s)$ occurs only in regions (values are at time $t$): replace each atom $\alpha$ mentioning it by $\bigvee_{(\pi',c')}(\pi'\wedge\alpha[z(s):=c'])$ over the pairs of $s_{z(s)}$. Repeat; each round brings in only variables before $z(s)$, so it ends with $\tilde s_y$.
\item \emph{Seed.} $C_{i_k}:=\big(\TRUE,\ F_k\wedge\bigwedge_{t\le H}\bigwedge_{y\in Y_k(t)}\delta_k\df{y=\tilde s_y}\big)$.
\item \emph{Sweep: quotient and projection.} For $k=1,\dots,p$, with current sets, $R_k:=A_P\wedge\bigwedge_{j\ne k}S_{i_j}$ (outside $N$: $\TRUE$), $X_k:=\hid{\ports(i_k)}$, $S_{i_k}:=\forall X_k(R_k\to G_P)$, the implementation set of $\rho_{\ports(i_k)}\big(P/\bigotimes_{j\ne k}(\TRUE,S_{i_j})\big)$: the quotient's guarantee is $S_P\vee\neg R_k=\neg A_P\vee G_P\vee\neg R_k\equiv R_k\to G_P$, as $\neg A_P\models\neg R_k$.
\item \emph{Assumptions: projection.} With the final sets, $A_{i_k}:=\exists X_k\,R_k$, the projected quotient's assumption. Output $(A_{i_k},S_{i_k})$.
\end{enumerate}
Steps 4, 8, 9 may use any equivalent formula of $\mathcal L$; the theorems do not depend on it, the witnesses do. The script uses Z3's \texttt{qe2} output, simplified and checked equivalent (not checked for one time per atom), and reads each witness from Z3's simplified $\psi$, which \texttt{DELEGATE\_TRACE=1} prints with $T$.

\section{Correctness}
\begin{theorem}[necessity]\label{thm:nec}
Under (Scope), for $\Phi$ from a topological ordering with (Pure): if an acceptable delegation exists, $\Phi$ is true.
\end{theorem}
\begin{proof}
$\Phi$ is true iff functions choosing each existential timed variable from the earlier ones make $A_P\to G_P$ hold for all universal values. By (K) each $i$ has a causal strategy $f_i$ with traces in $S_i$; choose $Y_k(t)$ as the output of $f_{i_k}$ at $t$. This uses only earlier values: an input of $i_k$ is an environment variable, in $E_j$ for the first $i_j$ ($j\le k$) reading it, or an output of some $i_j$, $j<k$; its values at times $\le t$ are quantified earlier, or are edge values that $f_{i_j}$ computes from such values. Complete the trace (other environment values arbitrary, components outside $N$ by their strategies): it lies in every $S_i$, so by (V1) it satisfies $A_P\to G_P$, which by (Scope), (Pure) mentions only quantified values.
\end{proof}
\begin{theorem}[soundness]\label{thm:sound}
If the ordering is nested and $\Phi$ is true, the seed of step 7 satisfies (V1), (V2), (R), (K), and each guarantee is a formula of $\mathcal L$ over $\ports(i_k)$.
\end{theorem}
\begin{proof}
\emph{Language.} $\df{y=\tilde s_y}\in\mathcal L$ by Lemma~\ref{lem:elim}(c) and step 6, and $\delta_k$ keeps this. \emph{Ports.} A timed variable of $\tilde s_y$, $y\in Y_k(t)$, precedes $y$, is not a terminal output of a later $i_j$, and lies in $E_j(s)$ or $Y_j(s)$ with $(s,j)$ not after $(t,k)$. If $j>k$ then $s<t$, so $H>0$ and $E_j(s)=\emptyset$, while $Y_j(s)$ was removed in step 6. So it is a port of some $i_j$, $j\le k$, at a time $\le t$, which $\delta_k$ maps into $\ports(i_k)$. \emph{(R), (K).} At time $t$, $i_k$ computes $Y_k(t)$ in prefix order from $\delta_k(\tilde s_y)$, from its inputs at times $\le t$ (Lemma~\ref{lem:carry}) and its outputs already computed, then sets $\sigma_k(v)(t):=v(t)$, other edge values $0$: a causal strategy meeting the guarantee. \emph{(V2)}: all assumptions are $\TRUE$. \emph{(V1).} If $x$ satisfies $A_P$ and every guarantee, it satisfies $F_1,\dots,F_{p-1}$, so $x(y)=\tilde s_y(x)$ by decoding, and $x(y)=s_y(x)$ by induction along the prefix (step 6 replaces $z(s)$, earlier, by $s_{z(s)}$, worth $x(z(s))$ at $x$). With $q_1,\dots,q_n$ the prefix, $\psi_n:=A_P\to G_P$ and $\psi_{l-1}$ the elimination of $q_l$ from $\psi_l$ ($\psi_0\equiv\TRUE$), induction gives $x\models\psi_l$: if $q_l$ is universal, $\psi_{l-1}\equiv\forall q_l\psi_l$; if existential, $\psi_{l-1}\models\psi_l[q_l:=s_{q_l}]$ and $s_{q_l}(x)=x(q_l)$. So $x\models G_P$.
\end{proof}
\begin{theorem}[sweep]\label{thm:sweep}
If the seed satisfies (V1), (R), (K), steps 8 and 9 output an acceptable delegation in $\mathcal L$ whose sets contain the seed's. It is maximal: no formula over $\ports(i_k)$ strictly weaker than $S_{i_k}$ satisfies (V1) with the other final sets, nor does any family weaker in every set and strictly in one. With the final sets, $(A_{i_k},S_{i_k})=\rho_{\ports(i_k)}\big(P/\bigotimes_{j\ne k}(\TRUE,S_{i_j})\big)=\rho_{\ports(i_k)}\big(P/\bigotimes_{j\ne k}C_{i_j}\big)$, the $C_{i_j}$ the output contracts.
\end{theorem}
\begin{proof}
Fix $k$ and the other sets; $X:=X_k$. For $S$ over $\ports(i_k)$: $R_k\wedge S\models G_P$ iff $S\models\forall X(R_k\to G_P)$ (if $x\models S$, so does every trace differing from $x$ only on $X$). So step 8 assigns the weakest $S$ satisfying (V1) with the current others; the previous one did, so sets only grow and inherit (R), (K). Let $S^*$ be the weakest for the final others: the final family satisfies (V1), so $S_{i_k}\models S^*$; after step $k$ the others only grew, so $R_k$ grew and $S^*\models S_{i_k}$. A family weaker everywhere satisfying (V1) still does after restoring all but one set, so each set is $\equiv$ the final one. Step 9: $A_P\wedge\bigwedge_jS_{i_j}\models R_k\models\exists XR_k=A_{i_k}$: (V2). If $x\models\neg A_{i_k}$, no trace agreeing with $x$ on $\ports(i_k)$ satisfies the final $R_k$, nor step 8's smaller one, so $x\models S_{i_k}$: saturated. As $S_{i_k}\equiv S^*$, the first equality is steps 8, 9; the second quotient has the same assumption and a guarantee $R_k\to(G_P\wedge\bigwedge_{j\ne k}A_{i_j})$, equivalent by Lemma~\ref{lem:adj} as the output is valid.
\end{proof}
So, on a nested ordering under (Scope), an acceptable delegation exists iff $\Phi$ is true, and steps 6 to 9 compute one: maximal, not greatest, and dependent on the orders of $T$, the prefix and the sweep.

\section{The relay, every step}\label{sec:relay}
\textbf{1.} $N=\{$Radio, Display$\}$ (Display writes $d$; Radio$\to$Display); (Scope) holds. \textbf{2.} Radio, Display is the topological ordering. $\priv_1=\{r,m\}\setminus\{m,d\}=\{r\}$, $\edge_1=\{m\}$; (E), (Cap) $1\le1$, (Pure), (Info) hold; $\sigma_1(r)=m$, $F_1=(m=r)$. \textbf{3.} $E_1=\{r\}$, $Y_1=\emptyset$, $E_2=\emptyset$, $Y_2=\{d\}$; $\Phi=\forall r\,\exists d\,(|r|\le1\to|d-r|\le\frac12)$. \textbf{4.} Pair $(0,2)$: the bounds of $d$ are $r\pm\frac12$; $T$ is $r$, $r-\frac12$, $r+\frac12$, $r+\frac32$, $r-\frac32$, $0$ ($r-\frac12+1$ and $r+\frac12-1$ are duplicates); $\psi[d:=r]\equiv\TRUE$, so $\pi_r\equiv\TRUE$, $s_d=r$, $\exists d\,\psi\equiv\TRUE$. Pair $(0,1)$: $\forall r\,\TRUE\equiv\TRUE$. \textbf{5.} $\Phi$ is true. \textbf{6.} Nothing to replace. \textbf{7.} $\df{d=s_d}\equiv(d=r)$ (as $\pi_r\equiv\TRUE$) and $\kappa_2(r)=m$, so $\delta_2$ gives $d=m$. Seed: Radio $(\TRUE,m=r)$, Display $(\TRUE,d=m)$: $D$ of Proposition~\ref{prop:nogreatest}.

\textbf{8, Radio.} $X_1=\{d\}$; $S_{\mathrm{Radio}}:=\forall d\,\big((|r|\le1\wedge d=m)\to|d-r|\le\frac12\big)\equiv\big(|r|\le1\to|m-r|\le\frac12\big)$, the premise holding only for $d=m$.
\textbf{8, Display.} $X_2=\{r\}$; $R_2\equiv|r|\le1\wedge|m-r|\le\frac12$; $S_{\mathrm{Display}}:=\forall r(R_2\to|d-r|\le\frac12)\equiv\neg\exists r\big(R_2\wedge(r>d+\frac12\vee r<d-\frac12)\big)$. Case $r>d+\frac12$: lower bounds $-1$, $m-\frac12$, $d+\frac12$ (strict), upper bounds $1$, $m+\frac12$; Fourier--Motzkin: $-1\le1$, $-1\le m+\frac12$, $m-\frac12\le1$, $m-\frac12\le m+\frac12$, $d+\frac12<1$, $d+\frac12<m+\frac12$, i.e.\ $|m|\le\frac32\wedge d<\frac12\wedge d<m$. Case $r<d-\frac12$: lower $-1$, $m-\frac12$; upper $1$, $m+\frac12$, $d-\frac12$ (strict); the new pairs give $d>-\frac12$, $d>m$: $|m|\le\frac32\wedge d>-\frac12\wedge d>m$. Negating,
\[S_{\mathrm{Display}}\equiv|m|\le\tfrac32\to\big((d\ge\tfrac12\vee d\ge m)\wedge(d\le-\tfrac12\vee d\le m)\big)\quad(d=m\text{ on }|m|\le\tfrac12).\]
\textbf{9.} $A_{\mathrm{Display}}:=\exists r(|r|\le1\wedge|m-r|\le\frac12)$: lower $-1$, $m-\frac12$, upper $1$, $m+\frac12$, so $|m|\le\frac32$. $A_{\mathrm{Radio}}:=\exists d(|r|\le1\wedge S_{\mathrm{Display}})\equiv|r|\le1$, as $d:=m$ satisfies $S_{\mathrm{Display}}$ for every $m$. Output: Radio $\big(|r|\le1,\ |r|\le1\to|m-r|\le\frac12\big)$, Display $\big(|m|\le\frac32,\ S_{\mathrm{Display}}\big)$. \textbf{False $\Phi$.} $H=2$, $A_P=\square_{[0,2]}|r|\le1$, $G_P=\square_{[0,2]}|d-r|\le\frac12\wedge\Diamond_{[0,2]}d\ge\frac14$, $\Phi$ is false ($r\equiv-1$ forces $d\le-\frac12$).


\section{Reproduction}\label{sec:repro}
\sloppy Install with \texttt{uv venv}, then \texttt{uv pip install -p} \texttt{.venv/bin/python} \texttt{z3-solver==5.3.1.0}; from the root run \texttt{.venv/bin/python scripts/delegate.py}. It prints the Z3 version, a block per instance and regression instances (tests with known verdicts), and exits $0$ iff all verdicts are as expected. ``receptive \& causal'' is $\forall\inp(0)\exists\outp(0)\cdots\forall\inp(H)\exists\outp(H)\,S_i$ (values at time $t$ of $\inp(i)$, $\outp(i)$): (K), read through functions picking each $\exists$-value from earlier ones (Lemma~\ref{lem:elim}(b)). A seed prints as \texttt{y == If(...)}, i.e.\ $\df{y=s_y}$; the sweep's \texttt{G} is $S_{i_k}$, a Boolean combination of atoms printed as Z3 simplifies it (order may vary with Z3's version). First 21 lines:\vspace{-3pt}
\begin{Verbatim}[fontsize=\fontsize{8}{8.8}\selectfont,breaklines=true,breakanywhere=true]
z3 version 5.3.1
==============================================================================
relay, H = 0
  A_P = And(r_0 <= 1, r_0 >= -1)
  G_P = And(d_0 + -1*r_0 <= 1/2, -1*d_0 + r_0 <= 1/2)
  N = ['Display', 'Radio']
  nested ordering: ['Radio', 'Display']
  priv(Radio) = ['r'], edge(Radio) = ['m'], forwarding {'r': 'm'}
  fresh inputs E_k: {'Radio': ['r'], 'Display': []}  terminal outputs Y_k: {'Radio': [], 'Display': ['d']}
  Phi (linearised sentence) = True
  contract Radio: A = true ; G = m_0 == r_0
  contract Display: A = true ; G = d_0 == m_0
  (V1) A_P & all G_i -> G_P : True
  Radio receptive & causal: True
  Display receptive & causal: True
  sweep to a maximal delegation (order ['Radio', 'Display'] ):
    Radio: A = And(r_0 <= 1, r_0 >= -1) ; G = Or(And(m_0 + -1*r_0 <= 1/2, -1*m_0 + r_0 <= 1/2),    Not(And(r_0 <= 1, r_0 >= -1)))
    Display: A = And(m_0 <= 3/2, m_0 >= -3/2) ; G = And(Or(0 <= -1*m_0 + d_0,        Not(m_0 >= -3/2),        1/2 <= d_0,        Not(m_0 <= 3/2)),     Or(Not(m_0 >= -3/2),        Not(m_0 <= 3/2),        0 <= m_0 + -1*d_0,        d_0 <= -1/2))
  (V1) after sweep: True
  (V2) after sweep: True
  receptive & causal after sweep (every component): True
\end{Verbatim}
\fussy
\vspace{-4pt}

\section{Scope and limits}\label{sec:scope}
\textbf{Forks.} Refutation still works: under (Scope), $\Phi^\circ$, built as in step 3 from any topological ordering but with $Y_k(t)$ all outputs of $i_k$ at $t$, is true if an acceptable delegation exists (proof of Theorem~\ref{thm:nec}, without (Pure)); so a refutation test decides $\Phi^\circ$ for every topological ordering. A true sentence proves nothing: let unconnected Left, Right map $s_1\mapsto a$, $s_2\mapsto b$, $A_P=\TRUE$, $G_P=|a-b|\le\frac12\wedge|b-s_2|\le\frac12$. For the ordering Right, Left, (Pure) holds and $\forall s_2\exists b\forall s_1\exists a\,G_P$ is true ($b=s_2$, $a=b$), but strategies $a=f(s_1)$, $b=g(s_2)$ would make $g$ bounded (fix $s_1$) and unbounded: none exists ((Cap) fails; Left, Right refutes). \emph{Remark (semi-decision, $H=0$).} Enumerate pairs ($B$, rational $\theta$): $f_i^\theta$ lists at most $B$ pairs (region, affine map of $\inp(i)$), each region a conjunction of at most $B$ atoms or negated atoms; the first region that holds gives the value (default $0$). Substitute in topological order and decide $\forall(A_P\to G_P)[\outp(i):=f_i^\theta(\inp(i))]$ by Lemma~\ref{lem:elim}; if true, output the acceptable $(\TRUE,\df{y=f_i^\theta})$. If an acceptable delegation exists, (K) and Lemma~\ref{lem:elim}(b) give such rational strategies (disjoint regions in disjunctive normal form), so it halts; else it runs forever; no bound on $B$ is known. The script runs one instance of this search, not the enumeration: on a fan-in graph Z3 decides the nonlinear $\exists\theta\forall(\text{environment})(A_P\to G_P)$ for one affine map per output on the whole space (60\,s timeout), then checks $\theta$ linearly.

\textbf{Other limits.} (Scope) makes $\Phi$ a sentence; if it fails, step 1 says ``does not apply''. Failure of (Pure) and unit delays are not covered: a delay needs $m(t+1)=r(t)$, not definable in $\mathcal L$ (with other values fixed, a formula gives a finite union of products in the $(m(t+1),r(t))$-plane, and a product inside the line is a point). \textbf{Cost.} Lemma~\ref{lem:elim}(a) multiplies the $n$ atoms by at most $n(n-1)/2+3n+1$ per elimination; step 2 may try $p!$ orders.
\end{document}
